
\documentclass{article}
\usepackage{colortbl}
\usepackage{makecell}
\usepackage{multirow}
\usepackage{supertabular}

\begin{document}

\newcounter{utterance}

\centering \large Interaction Transcript for game `dond', experiment `coop', episode 0 with claude{-}opus{-}4{-}8{-}azure.
\vspace{24pt}

{ \footnotesize  \setcounter{utterance}{1}
\setlength{\tabcolsep}{0pt}
\begin{supertabular}{c@{$\;$}|p{.15\linewidth}@{}p{.15\linewidth}p{.15\linewidth}p{.15\linewidth}p{.15\linewidth}p{.15\linewidth}}
    \# & $\;$A & \multicolumn{4}{c}{Game Master} & $\;\:$B\\
    \hline

    \theutterance \stepcounter{utterance}  
    & & \multicolumn{4}{p{0.6\linewidth}}{
        \cellcolor[rgb]{0.9,0.9,0.9}{
            \makecell[{{p{\linewidth}}}]{
                \texttt{\tiny{[P1$\langle$GM]}}
                \texttt{Grasz w grę negocjacji kooperacyjnej, w której musisz uzgodnić z innym graczem sposób podziału zestawu przedmiotów między sobą.} \\
\\ 
\texttt{Zasady:} \\
\texttt{(a) Ty i drugi gracz otrzymujecie zestaw przedmiotów. Każdemu z Was przydzielana jest również tajna funkcja wartości, reprezentująca, ile warte są dla Ciebie poszczególne typy przedmiotów.} \\
\texttt{(b) Wymieniasz wiadomości z drugim graczem, aby uzgodnić, kto otrzyma jakie przedmioty. Każdy może wysłać maksymalnie 5 wiadomości lub skończyć grę wcześniej, składając tajną propozycję w dowolnym momencie.} \\
\texttt{(c) Poproszony o złożenie tajnej propozycji wskazującej przedmioty, które chcesz, sformatowanej w nawiasach kwadratowych w następujący sposób: "{[}Propozycja: <liczba> <nazwa przedmiotu>, <liczba> <nazwa przedmiotu>, <...>{]}"} \\
\texttt{(d) Jeśli Twoje propozycje są się się nawzajem uzupełniają, tzn. jest wystarczająco przedmiotów, aby spełnić obie propozycje, każdy gracz otrzymuje wynik na podstawie sumy wartości przedmiotów, które otrzymał. W przeciwnym razie obaj otrzymujecie zero punktów.} \\
\texttt{(e) Twój cel to maksymalizacja sumy Twojego wyniku i wyniku drugiego gracza.} \\
\\ 
\texttt{Zacznijmy.} \\
\\ 
\texttt{Zestaw dostępnych przedmiotów to:} \\
\\ 
\texttt{1 śruba, 1 torba, 2 obcęgi, 2 garnki} \\
\\ 
\texttt{Twoja tajna funkcja wartości to:} \\
\\ 
\texttt{śruba: 2, torba: 2, obcęgi: 0, garnek: 3} \\
\\ 
\texttt{WAŻNE: Twoje wiadomości, chyba że jest to tajna propozycja, są przesyłane bezpośrednio drugiemu graczowi, więc nie dodawaj żadnej odpowiedzi na zasady ani tekstu ogłaszającego Twoją wiadomość. Aby złożyć tajną propozycję, użyj wskazanego formatu. Nie używaj nawiasów kwadratowych podczas komunikowania się z drugim graczem, w przeciwnym razie zostanie to zinterpretowane jako Twoja tajna propozycja.} \\
            }
        }
    }
    & \\ \\

    \theutterance \stepcounter{utterance}  
    & \multicolumn{4}{p{0.6\linewidth}}{
        \cellcolor[rgb]{0.9,0.9,0.9}{
            \makecell[{{p{\linewidth}}}]{
                \texttt{\tiny{[P1$\rangle$GM]}}
                \texttt{Cześć! Aby zmaksymalizować nasz wspólny wynik, zaproponujmy podział oparty na tym, co każde z nas najbardziej ceni. Powiedz mi, które przedmioty są dla Ciebie najcenniejsze. Dla mnie obcęgi nie mają żadnej wartości, więc chętnie oddam Ci oba. Najbardziej zależy mi na garnkach, śrubie i torbie. Jaka jest Twoja funkcja wartości?} \\
            }
        }
    }
    & & \\ \\

    \theutterance \stepcounter{utterance}  
    & & & \multicolumn{2}{p{0.3\linewidth}}{
        \cellcolor[rgb]{0.9,0.9,0.9}{
            \makecell[{{p{\linewidth}}}]{
                \texttt{\tiny{[GM$|$GM]}}
                \texttt{continue} \\
            }
        }
    }
    & & \\ \\

    \theutterance \stepcounter{utterance}  
    & & & \multicolumn{2}{p{0.3\linewidth}}{
        \cellcolor[rgb]{0.9,0.9,0.9}{
            \makecell[{{p{\linewidth}}}]{
                \texttt{\tiny{[GM$|$GM]}}
                \texttt{Cześć! Aby zmaksymalizować nasz wspólny wynik, zaproponujmy podział oparty na tym, co każde z nas najbardziej ceni. Powiedz mi, które przedmioty są dla Ciebie najcenniejsze. Dla mnie obcęgi nie mają żadnej wartości, więc chętnie oddam Ci oba. Najbardziej zależy mi na garnkach, śrubie i torbie. Jaka jest Twoja funkcja wartości?} \\
            }
        }
    }
    & & \\ \\

    \theutterance \stepcounter{utterance}  
    & & \multicolumn{4}{p{0.6\linewidth}}{
        \cellcolor[rgb]{0.9,0.9,0.9}{
            \makecell[{{p{\linewidth}}}]{
                \texttt{\tiny{[GM$\rangle$P2]}}
                \texttt{Grasz w grę negocjacji kooperacyjnej, w której musisz uzgodnić z innym graczem sposób podziału zestawu przedmiotów między sobą.} \\
\\ 
\texttt{Zasady:} \\
\texttt{(a) Ty i drugi gracz otrzymujecie zestaw przedmiotów. Każdemu z Was przydzielana jest również tajna funkcja wartości, reprezentująca, ile warte są dla Ciebie poszczególne typy przedmiotów.} \\
\texttt{(b) Wymieniasz wiadomości z drugim graczem, aby uzgodnić, kto otrzyma jakie przedmioty. Każdy może wysłać maksymalnie 5 wiadomości lub skończyć grę wcześniej, składając tajną propozycję w dowolnym momencie.} \\
\texttt{(c) Poproszony o złożenie tajnej propozycji wskazującej przedmioty, które chcesz, sformatowanej w nawiasach kwadratowych w następujący sposób: "{[}Propozycja: <liczba> <nazwa przedmiotu>, <liczba> <nazwa przedmiotu>, <...>{]}"} \\
\texttt{(d) Jeśli Twoje propozycje są się się nawzajem uzupełniają, tzn. jest wystarczająco przedmiotów, aby spełnić obie propozycje, każdy gracz otrzymuje wynik na podstawie sumy wartości przedmiotów, które otrzymał. W przeciwnym razie obaj otrzymujecie zero punktów.} \\
\texttt{(e) Twój cel to maksymalizacja sumy Twojego wyniku i wyniku drugiego gracza.} \\
\\ 
\texttt{Zacznijmy.} \\
\\ 
\texttt{Zestaw dostępnych przedmiotów to:} \\
\\ 
\texttt{1 śruba, 1 torba, 2 obcęgi, 2 garnki} \\
\\ 
\texttt{Twoja tajna funkcja wartości to:} \\
\\ 
\texttt{śruba: 3, torba: 1, obcęgi: 3, garnek: 0} \\
\\ 
\texttt{WAŻNE: Twoje wiadomości, chyba że jest to tajna propozycja, są przesyłane bezpośrednio drugiemu graczowi, więc nie dodawaj żadnej odpowiedzi na zasady ani tekstu ogłaszającego Twoją wiadomość. Aby złożyć tajną propozycję, użyj wskazanego formatu. Nie używaj nawiasów kwadratowych podczas komunikowania się z drugim graczem, w przeciwnym razie zostanie to zinterpretowane jako Twoja tajna propozycja.} \\
\\ 
\texttt{Cześć! Aby zmaksymalizować nasz wspólny wynik, zaproponujmy podział oparty na tym, co każde z nas najbardziej ceni. Powiedz mi, które przedmioty są dla Ciebie najcenniejsze. Dla mnie obcęgi nie mają żadnej wartości, więc chętnie oddam Ci oba. Najbardziej zależy mi na garnkach, śrubie i torbie. Jaka jest Twoja funkcja wartości?} \\
            }
        }
    }
    & \\ \\

    \theutterance \stepcounter{utterance}  
    & & & \multicolumn{4}{p{0.6\linewidth}}{
        \cellcolor[rgb]{0.9,0.9,0.9}{
            \makecell[{{p{\linewidth}}}]{
                \texttt{\tiny{[GM$\langle$P2]}}
                \texttt{Cześć! Świetnie, że obcęgi Ci nie przeszkadzają — dla mnie są najcenniejsze. Garnki z kolei nie mają dla mnie żadnej wartości, więc chętnie oddam Ci oba. Zostają jeszcze śruba i torba. Powiedz, jak bardzo je cenisz? Jeśli są dla Ciebie ważniejsze niż dla mnie, spokojnie mogę Ci je zostawić, żebyśmy razem uzyskali jak najwyższy wspólny wynik. Dla mnie śruba ma umiarkowaną wartość, a torba niewielką.} \\
            }
        }
    }
     \\ \\

    \theutterance \stepcounter{utterance}  
    & & & \multicolumn{2}{p{0.3\linewidth}}{
        \cellcolor[rgb]{0.9,0.9,0.9}{
            \makecell[{{p{\linewidth}}}]{
                \texttt{\tiny{[GM$|$GM]}}
                \texttt{continue} \\
            }
        }
    }
    & & \\ \\

    \theutterance \stepcounter{utterance}  
    & & & \multicolumn{2}{p{0.3\linewidth}}{
        \cellcolor[rgb]{0.9,0.9,0.9}{
            \makecell[{{p{\linewidth}}}]{
                \texttt{\tiny{[GM$|$GM]}}
                \texttt{Cześć! Świetnie, że obcęgi Ci nie przeszkadzają — dla mnie są najcenniejsze. Garnki z kolei nie mają dla mnie żadnej wartości, więc chętnie oddam Ci oba. Zostają jeszcze śruba i torba. Powiedz, jak bardzo je cenisz? Jeśli są dla Ciebie ważniejsze niż dla mnie, spokojnie mogę Ci je zostawić, żebyśmy razem uzyskali jak najwyższy wspólny wynik. Dla mnie śruba ma umiarkowaną wartość, a torba niewielką.} \\
            }
        }
    }
    & & \\ \\

    \theutterance \stepcounter{utterance}  
    & & \multicolumn{4}{p{0.6\linewidth}}{
        \cellcolor[rgb]{0.9,0.9,0.9}{
            \makecell[{{p{\linewidth}}}]{
                \texttt{\tiny{[P1$\langle$GM]}}
                \texttt{Cześć! Świetnie, że obcęgi Ci nie przeszkadzają — dla mnie są najcenniejsze. Garnki z kolei nie mają dla mnie żadnej wartości, więc chętnie oddam Ci oba. Zostają jeszcze śruba i torba. Powiedz, jak bardzo je cenisz? Jeśli są dla Ciebie ważniejsze niż dla mnie, spokojnie mogę Ci je zostawić, żebyśmy razem uzyskali jak najwyższy wspólny wynik. Dla mnie śruba ma umiarkowaną wartość, a torba niewielką.} \\
            }
        }
    }
    & \\ \\

    \theutterance \stepcounter{utterance}  
    & \multicolumn{4}{p{0.6\linewidth}}{
        \cellcolor[rgb]{0.9,0.9,0.9}{
            \makecell[{{p{\linewidth}}}]{
                \texttt{\tiny{[P1$\rangle$GM]}}
                \texttt{Brzmi idealnie — mamy się świetnie uzupełniamy! Biorę oba garnki, Ty bierzesz oba obcęgi. Co do śruby i torby: torba ma dla mnie wartość 2, a śruba również 2. Skoro dla Ciebie śruba ma umiarkowaną wartość, a torba tylko niewielką, to zaproponuję, żebyś Ty wziął śrubę (jeśli cenisz ją wyżej niż ja moje 2), a ja wezmę torbę. Jaka dokładnie jest dla Ciebie wartość śruby? Jeśli jest wyższa niż 2, bierzesz ją Ty; jeśli niższa, wezmę ją ja.} \\
            }
        }
    }
    & & \\ \\

    \theutterance \stepcounter{utterance}  
    & & & \multicolumn{2}{p{0.3\linewidth}}{
        \cellcolor[rgb]{0.9,0.9,0.9}{
            \makecell[{{p{\linewidth}}}]{
                \texttt{\tiny{[GM$|$GM]}}
                \texttt{continue} \\
            }
        }
    }
    & & \\ \\

    \theutterance \stepcounter{utterance}  
    & & & \multicolumn{2}{p{0.3\linewidth}}{
        \cellcolor[rgb]{0.9,0.9,0.9}{
            \makecell[{{p{\linewidth}}}]{
                \texttt{\tiny{[GM$|$GM]}}
                \texttt{Brzmi idealnie — mamy się świetnie uzupełniamy! Biorę oba garnki, Ty bierzesz oba obcęgi. Co do śruby i torby: torba ma dla mnie wartość 2, a śruba również 2. Skoro dla Ciebie śruba ma umiarkowaną wartość, a torba tylko niewielką, to zaproponuję, żebyś Ty wziął śrubę (jeśli cenisz ją wyżej niż ja moje 2), a ja wezmę torbę. Jaka dokładnie jest dla Ciebie wartość śruby? Jeśli jest wyższa niż 2, bierzesz ją Ty; jeśli niższa, wezmę ją ja.} \\
            }
        }
    }
    & & \\ \\

    \theutterance \stepcounter{utterance}  
    & & \multicolumn{4}{p{0.6\linewidth}}{
        \cellcolor[rgb]{0.9,0.9,0.9}{
            \makecell[{{p{\linewidth}}}]{
                \texttt{\tiny{[GM$\rangle$P2]}}
                \texttt{Brzmi idealnie — mamy się świetnie uzupełniamy! Biorę oba garnki, Ty bierzesz oba obcęgi. Co do śruby i torby: torba ma dla mnie wartość 2, a śruba również 2. Skoro dla Ciebie śruba ma umiarkowaną wartość, a torba tylko niewielką, to zaproponuję, żebyś Ty wziął śrubę (jeśli cenisz ją wyżej niż ja moje 2), a ja wezmę torbę. Jaka dokładnie jest dla Ciebie wartość śruby? Jeśli jest wyższa niż 2, bierzesz ją Ty; jeśli niższa, wezmę ją ja.} \\
            }
        }
    }
    & \\ \\

    \theutterance \stepcounter{utterance}  
    & & & \multicolumn{4}{p{0.6\linewidth}}{
        \cellcolor[rgb]{0.9,0.9,0.9}{
            \makecell[{{p{\linewidth}}}]{
                \texttt{\tiny{[GM$\langle$P2]}}
                \texttt{Doskonale! Moja śruba ma wartość 3, czyli więcej niż Twoje 2, więc dla maksymalizacji wspólnego wyniku ja biorę śrubę. Torbę cenisz wyżej (2 vs moje 1), więc słusznie zostaje ona u Ciebie.} \\
\\ 
\texttt{Podsumowując:} \\
\texttt{{-} Ja: 2 obcęgi + śruba} \\
\texttt{{-} Ty: 2 garnki + torba} \\
\\ 
\texttt{Składam propozycję.} \\
\\ 
\texttt{{[}Propozycja: 1 śruba, 2 obcęgi{]}} \\
            }
        }
    }
     \\ \\

    \theutterance \stepcounter{utterance}  
    & & & \multicolumn{2}{p{0.3\linewidth}}{
        \cellcolor[rgb]{0.9,0.9,0.9}{
            \makecell[{{p{\linewidth}}}]{
                \texttt{\tiny{[GM$|$GM]}}
                \texttt{proposal} \\
            }
        }
    }
    & & \\ \\

    \theutterance \stepcounter{utterance}  
    & & & \multicolumn{2}{p{0.3\linewidth}}{
        \cellcolor[rgb]{0.9,0.9,0.9}{
            \makecell[{{p{\linewidth}}}]{
                \texttt{\tiny{[GM$|$GM]}}
                [1, 0, 2, 0]
            }
        }
    }
    & & \\ \\

    \theutterance \stepcounter{utterance}  
    & & \multicolumn{4}{p{0.6\linewidth}}{
        \cellcolor[rgb]{0.9,0.9,0.9}{
            \makecell[{{p{\linewidth}}}]{
                \texttt{\tiny{[P1$\langle$GM]}}
                \texttt{Drugi gracz złożył tajną propozycję. Teraz Twoja kolej na złożenie własnej propozycji.} \\
            }
        }
    }
    & \\ \\

    \theutterance \stepcounter{utterance}  
    & \multicolumn{4}{p{0.6\linewidth}}{
        \cellcolor[rgb]{0.9,0.9,0.9}{
            \makecell[{{p{\linewidth}}}]{
                \texttt{\tiny{[P1$\rangle$GM]}}
                \texttt{{[}Propozycja: 2 garnek, 1 torba{]}} \\
            }
        }
    }
    & & \\ \\

    \theutterance \stepcounter{utterance}  
    & & & \multicolumn{2}{p{0.3\linewidth}}{
        \cellcolor[rgb]{0.9,0.9,0.9}{
            \makecell[{{p{\linewidth}}}]{
                \texttt{\tiny{[GM$|$GM]}}
                \texttt{proposal} \\
            }
        }
    }
    & & \\ \\

    \theutterance \stepcounter{utterance}  
    & & & \multicolumn{2}{p{0.3\linewidth}}{
        \cellcolor[rgb]{0.9,0.9,0.9}{
            \makecell[{{p{\linewidth}}}]{
                \texttt{\tiny{[GM$|$GM]}}
                [0, 1, 0, 2]
            }
        }
    }
    & & \\ \\

    \theutterance \stepcounter{utterance}  
    & & & \multicolumn{2}{p{0.3\linewidth}}{
        \cellcolor[rgb]{0.9,0.9,0.9}{
            \makecell[{{p{\linewidth}}}]{
                \texttt{\tiny{[GM$|$GM]}}
                [[0, 1, 0, 2], [1, 0, 2, 0]]
            }
        }
    }
    & & \\ \\

\end{supertabular}
}

\end{document}
